\hypertarget{cpython-memory-allocator-pymalloc-generational-garbage-collection}{%
\chapter{CPython Memory Allocator (PyMalloc) \& Generational Garbage
Collection}\label{cpython-memory-allocator-pymalloc-generational-garbage-collection}}

\hypertarget{cpythons-memory-allocation-engine-pymalloc}{%
\subsection{23.1 CPython's Memory Allocation Engine
(PyMalloc)}\label{cpythons-memory-allocation-engine-pymalloc}}

For large allocations (greater than 512 bytes), CPython forwards the
request directly to the system's standard C library
\passthrough{\lstinline!malloc()!}. However, for small objects
(\(\le 512\) bytes)which represent the vast majority of Python
allocationsstandard operating system allocators introduce high
fragmentation and locking overhead. To resolve this, CPython implements
a custom small-object allocator called \textbf{PyMalloc}.

\hypertarget{pymalloc-memory-hierarchy}{%
\subsubsection{1. PyMalloc Memory
Hierarchy}\label{pymalloc-memory-hierarchy}}

PyMalloc structures memory into three layers:

\begin{lstlisting}
+-------------------------------------------------------------+
|                     ARENA (256 KB)                          |
|  Aligned to 256 KB boundaries in virtual memory.            |
|  Contains exactly 64 Pools.                                 |
|                                                             |
|   +-------------------+ +-------------------+               |
|   |    POOL (4 KB)    | |    POOL (4 KB)    |  ...          |
|   |  Allocates one    | |  Allocates one    |               |
|   |  size-class only. | |  size-class only. |               |
|   |  +--------------+ | |  +--------------+ |               |
|   |  | BLOCK (8B-512B)| |  | BLOCK (8B-512B)| |               |
|   |  +--------------+ | |  +--------------+ |               |
|   +-------------------+ +-------------------+               |
+-------------------------------------------------------------+
\end{lstlisting}

\begin{enumerate}
\def\labelenumi{\arabic{enumi}.}
\tightlist
\item
  \textbf{Arenas}: Contiguous virtual memory blocks of 256 KB, aligned
  to 256 KB boundaries. Arenas manage virtual memory allocations from
  the OS and contain exactly 64 pools.
\item
  \textbf{Pools}: Blocks of 4 KB (matching the OS virtual page size). A
  pool is dedicated to a single \textbf{size class} (e.g., all blocks in
  a pool are 32 bytes, or all are 64 bytes).
\item
  \textbf{Blocks}: The actual chunks of memory returned to the
  interpreter. Blocks range from 8 bytes to 512 bytes, aligned to 8-byte
  steps (giving 64 distinct size classes).
\end{enumerate}

\hypertarget{size-class-formula}{%
\subsubsection{2. Size Class Formula}\label{size-class-formula}}

The size class for an allocation request of size \(S\) is calculated by:
\[\text{Class Index} = \left\lceil \frac{S}{8} \right\rceil - 1\] This
structured layout allows PyMalloc to bypass the general-purpose
allocator heap search and perform near-instantaneous \(\mathcal{O}(1)\)
allocation and free operations using thread-local lookups.

\begin{center}\rule{0.5\linewidth}{0.5pt}\end{center}

\hypertarget{generational-garbage-collection-cycle-detection}{%
\subsection{23.2 Generational Garbage Collection Cycle
Detection}\label{generational-garbage-collection-cycle-detection}}

While reference counting manages the vast majority of object lifecycles,
it cannot identify self-referencing cyclic loops (e.g.,
\passthrough{\lstinline!x = []; x.append(x)!}). To reclaim cyclic memory
leaks, CPython runs a cyclic Garbage Collector (GC) as a background
system.

\hypertarget{c-level-gc-header-layout}{%
\subsubsection{1. C-level GC Header
Layout}\label{c-level-gc-header-layout}}

Every object tracked by the GC has an extra header prefix in memory
before its normal \passthrough{\lstinline!PyObject!} structure. The GC
cast looks up objects using the \passthrough{\lstinline!PyGC\_Head!}
struct:

\begin{lstlisting}[language=C]
typedef union _gc_head {
    struct {
        union _gc_head *gc_next;    /* Pointer to next object in the generation list */
        union _gc_head *gc_prev;    /* Pointer to previous object in the generation list */
        Py_ssize_t gc_refs;         /* GC reference count copy used during cycle checks */
    } gc;
    double dummy;                   /* Forces alignment boundary adjustments */
} PyGC_Head;
\end{lstlisting}

\hypertarget{the-three-generations}{%
\subsubsection{2. The Three Generations}\label{the-three-generations}}

The GC divides tracked objects into three generations based on survival
age: * \textbf{Generation 0 (Gen 0)}: Where new objects are registered.
Collected frequently. * \textbf{Generation 1 (Gen 1)}: Objects that
survived a Gen 0 collection pass. * \textbf{Generation 2 (Gen 2)}:
Long-lived objects. Collected least frequently.

Collection is triggered when the number of allocations minus
deallocations in a generation exceeds a configured threshold:
\[\text{Allocations} - \text{Deallocations} > \text{Threshold}\]

\hypertarget{the-cycle-detection-algorithm}{%
\subsubsection{3. The Cycle Detection
Algorithm}\label{the-cycle-detection-algorithm}}

To find and break reference cycles without locking the entire runtime
heap, CPython executes the following steps: 1. \textbf{Copy Reference
Counts}: The GC copy-assigns the reference count of every object in the
collection generation to its \passthrough{\lstinline!gc\_refs!} field in
the \passthrough{\lstinline!PyGC\_Head!} struct. 2. \textbf{Subtract
Internal References}: For each object, the GC traverses its linked
references (via \passthrough{\lstinline!tp\_traverse!} slots). If a
target object is also in the collection generation, the GC decrements
the target's \passthrough{\lstinline!gc\_refs!} count. 3.
\textbf{Isolate Garbage}: * If an object's
\passthrough{\lstinline!gc\_refs!} drops to 0, it means the object is
only reachable via references \emph{within} the generation cycle
(candidates for garbage). * If \passthrough{\lstinline!gc\_refs > 0!},
the object is reachable from outside the generation. The object and all
objects it references are marked as reachable and moved back to the main
list. 4. \textbf{Deallocate Cycles}: Any remaining objects with
\passthrough{\lstinline!gc\_refs == 0!} are confirmed unreachable. The
GC breaks the reference cycle by clearing references and deallocating
the memory.

\begin{center}\rule{0.5\linewidth}{0.5pt}\end{center}
